\input{../tex-inputs/latex-leseansicht-vorspann}

\section[Gustav Schwarzkopf an Arthur Schnitzler, 21. 8. 1913]{L04262 Gustav Schwarzkopf an Arthur Schnitzler, 21. 8. 1913}
\nopagebreak\mylabel{L04262v}
\rehead{ }\normalsize\beginnumbering\briefempfaengerindex{Schnitzler, Arthur@\textsc{Schnitzler, Arthur}!zzzSchwarzkopf, Gustav@\emph{von Gustav Schwarzkopf}!1913-08-211@{21. 8. 1913}|(be}
\toendnotes[C]{\smallbreak\pagebreak[2]}
\correspDesc{Versand  durch Gustav Schwarzkopf am 21. 8. 1913 in Wien
\newline{}Übermittlung  am 22. 8. 1913 in Wien
\newline{}Erhalt  durch Arthur Schnitzler am 23. 8. 1913 in Venedig}\toendnotes[C]{\smallbreak}
\Standort{CUL, Schnitzler, B 96a.}
\physDesc{Postkarte, 680 Zeichen
\newline{}Handschrift: schwarze Tinte, deutsche Kurrent
\newline{}Versand: 1) Stempel: »\nobreak{}\oindex{I., Innere Stadt@\textbf{I., Innere Stadt}, \emph{Verwaltungsgebiet}|pwk}1/1 Wien 8, 22. VI\textcolor{gray}{II}. 13, IX\nobreak{}«.   2) Stempel: »\nobreak{}\oindex{Venedig@\textbf{Venedig}|pwk}Venezia Centro, 23. 8. 13, 12M\nobreak{}«. 
\newline{}Schnitzler: mit Bleistift Vermerk: »\textsc{Schwarzk}{[}opf{]}« }\toendnotes[C]{\smallbreak}\pstart{}{\pb}I. Tiefer Graben 17\oindex{Wien@\textbf{Wien}!I., Innere Stadt@\textbf{I., Innere Stadt}!Tiefer Graben 17@\textbf{Tiefer Graben 17}, \emph{Wohngebäude}|pw}\pend{}\pstart{}II. 6\pend{}{\bigskip}\pstart{}Herrn\pend{}\pstart{}\textsc{D\textsuperscript{r} Arthur Schnitzler}\pend{}\pstart{}\textsc{Venedig\oindex{Venedig@\textbf{Venedig}|pw}}\pend{}\pstart{}\textsc{Grand Hotel\oindex{Grand Hotel Venezia@\textbf{Grand Hotel Venezia}, \emph{Hotel}|pw}.}\pend{}{\bigskip}\vspace{1em}
\pstart
           \raggedleft{}{\pb}21. VIII. 13\pend
           \vspace{0.5em}
\pstart
           Lieber Arthur, es geht mir eigentlich wieder ganz gut, oder –
               genauer ausgedrückt – mein Befinden iſt wieder ſo, wie es vor dem plötzlichen Unfall
               war. We{\geminationn} nur das Waſſer ein bischen freundlicher wäre.
               Aber oft iſt es ſo, daß es nicht einmal für einen Spaziergang in den Volksgarten\oindex{Dänemark@\textbf{Dänemark}|pw} reicht. Einen der wenigen Abende, die \uline{nicht} verregnet waren habe\strikeout{n} ich mit \textsc{Paul Marx\pwindex{Marx, Paul 21.\,7.\,1879 Wien – 30.\,10.\,1956 ebd.@\textsc{Marx, Paul} (21.\,7.\,1879 Wien – 30.\,10.\,1956 ebd.), \emph{Regisseur, Schauspieler}|pw}} im Prater\oindex{Wien@\textbf{Wien}!II., Leopoldstadt@\textbf{II., Leopoldstadt}!Prater@\textbf{Prater}, \emph{Park}|pw} verbracht. Er fährt am 23.
               nach München\oindex{München@\textbf{München}|pw} und läßt herzlichſt grüßen.\pend
           
\pstart
           Ich danke Ihnen für Ihre Erkundigung, wünſche Ihnen, ob Sie Ihr Weg nun nach dem Engadin\oindex{Dänemark@\textbf{Dänemark}|pw} oder nach Oberitalien\oindex{Italien@\textbf{Italien}|pw} führt, gutes Wetter!\pend
           
\pstart
           {\pb}und grüße Ihre Frau\pwindex{Schnitzler, Olga 17.\,1.\,1882 Wien – 13.\,1.\,1970 Lugano@\textsc{Schnitzler, Olga} (17.\,1.\,1882 Wien – 13.\,1.\,1970 Lugano), \emph{Schauspielerin, Sängerin}|pwv}{\\[\baselineskip]}und Sie herzlichſt{\\[\baselineskip]}Ihr{\\[\baselineskip]}\spacefill\mbox{Gustav}\pend
           \leftskip=0em{}\selectlanguage{ngerman}\endnumbering\briefempfaengerindex{Schnitzler, Arthur@\textsc{Schnitzler, Arthur}!zzzSchwarzkopf, Gustav@\emph{von Gustav Schwarzkopf}!1913-08-211@{21. 8. 1913}|)be}\mylabel{L04262h}
\begin{anhang}
\end{anhang}\newcommand{\dateiname}{L04262}\newcommand{\titel}{Gustav Schwarzkopf an Arthur Schnitzler, 21. 8. 1913}\newcommand{\editorInnen}{Herausgegeben von Jahnke, SelmaMüller, Martin Anton}\input{../tex-inputs/latex-leseansicht-abspann}
